\chapter{Conclusion and Future Work}
\label{ch:conclusion}

This thesis built, validated, and optimised a from-scratch CUDA/OptiX volumetric path tracer for
Gaussian kernel-mixture volumes---a high-performance re-implementation of \emph{Don't Splat Your
Gaussians}~\cite{DSYG}---one that, at equal quality, renders faster than the Mitsuba reference.

\section{Summary of contributions}
\label{sec:contributions}

The central contribution is a rendering architecture (\Cref{ch:architecture})---suggested by Condor,
the lead author of the reference (personal communication)---with two
complementary halves: \emph{single-trace any-hit collection}, which gathers all of a ray's primitive
entries in one acceleration-structure traversal (back-face-culled to entries, exits computed
analytically), and \emph{analog-decomposition scatter sampling}, in which each primitive draws a
closed-form free flight and the nearest---the argmin---is the scattering event. Together they replace
the segment-marching and root-finding that the reference uses for the scatter
distance. This rests on two supporting pillars: a \emph{differential-validation methodology}
(\Cref{ch:validation}) that establishes correctness against Mitsuba and, where a closed form exists,
against analytic ground truth---and that doubles as the unbiasedness evidence for the new sampler; and a
systematic \emph{performance-engineering study} (\Cref{ch:optimization}) that contributed a second,
scene-dependent win (volumetric product-RIS direct lighting---$1.48\times$ at equal quality under
environment lighting, default off) and a characterised ledger of negative results. At equal
quality the renderer renders faster than the reference on the showcase, and free of the fireflies the
reference exhibits there. The advantage is not confined to frame time: at equal
conditions the reference's faster direct-lighting estimator is energy-biased on dense media while this
renderer's is not (\Cref{sec:results-firefly}), per-asset cap calibration puts the renderer's peak
device memory below the reference's on the showcase (\Cref{sec:results-memory}), and its
ahead-of-time-compiled pipeline starts in roughly half the reference's warm start
(\Cref{sec:results-startup}). Finally, a cross-architecture probe (\Cref{sec:ser}) shows the
hand-written megakernel exploits Shader Execution Reordering on Ada hardware for a further
\num{1.12}--\num{1.68}$\times$ at equal quality, image-identical---a divergence lever the reference's
compiled pipeline cannot reach. Measured end-to-end on that hardware against the same reference build, the
equal-quality advantage over the unbiased reference grows from the \num{59}$\times$ of the \num{3090} to
\num{110}$\times$ with reordering, and to \num{151}$\times$ with the renderer's best unbiased
configuration (\Cref{sec:ser})---Ada figures that corroborate, and do not replace, the pinned
\num{59}$\times$.

\section{Limitations}
\label{sec:limitations}

The renderer is deliberately scoped, and several limitations bound its current reach:

\begin{itemize}
  \item \textbf{No emission.} It models absorption and scattering but not emissive media, so combustion
  assets render as dark plumes; the assets carry no emission data either.
  \item \textbf{Gaussian kernel only.} The closed-form optical depth (\Cref{sec:analytic-od}) is
  Gaussian-specific, so other kernel families---e.g.\ the compactly-supported Epanechnikov kernel---are
  unsupported; each would need its own closed-form integral, which the architecture could accommodate
  but does not implement. Relatedly, assets fit with the Gabor-noise extension render only their coarse
  Gaussian base; high-frequency Gabor detail is unsupported.
  \item \textbf{Static, single-frame.} There is no animation support.
  \item \textbf{Compile-time, asset-dependent caps.} The per-ray buffers (\Cref{sec:gpu-impl}) are sized
  at compile time; a substantially denser scene requires re-sizing and a rebuild. An offline estimator
  now predicts the required caps from an asset's geometry---the \num{25600}-Gaussian bunny, for one,
  needs markedly larger caps than the \num{652}-primitive cloud---so the sizing is determined by
  measurement rather than assumption, but the caps remain static, and the runtime overflow counter only
  \emph{detects} an undersized cap rather than fixing it.
  \item \textbf{Ampere hardware.} The one lever that directly targets the renderer's ray divergence,
  OptiX Shader Execution Reordering~\cite{NvidiaSER2022}, is Ada-only and unavailable on the RTX~3090
  used throughout---though its effect is measured directly on an Ada GPU in \Cref{sec:ser}
  (\num{1.12}--\num{1.68}$\times$, image-identical).
  \item \textbf{Small systematic residuals.} A handful of sub-$0.2\%$ systematic differences from the
  reference are documented rather than resolved: a ${+}2\times10^{-4}$ difference at heavily-overlapped
  scatter vertices (candidate causes: the argmin free-flight distribution under overlap, or next-event
  shadow-ray transmittance from a vertex inside overlapping primitives; \Cref{ch:validation}); a ${+}1.8\times10^{-4}$ brightness systematic in low-density interiors; and a
  small channel-dependent residual in coloured-albedo rendering. Each is bounded well under a part in
  $500$ and confined to a specific regime.
\end{itemize}

\section{Future work}
\label{sec:future-work}

The most direct continuations are natural extensions of the kernel-mixture representation, several
suggested by Condor: a \emph{unified emissive/absorptive kernel} model that exploits closed-form
integrals for accurate product sampling of light sources and BSDFs; \emph{SGGX and microflake} models
to represent surface-like materials within the same primitive framework; \emph{parametric kernels} for
more compact and accurate radiance fields; and a bridge from \emph{Lagrangian fluid simulation} directly
to kernel mixtures, bypassing voxelisation and kernel regression.

Three engineering directions follow from this thesis specifically. First, \emph{differentiable and
inverse rendering}: re-expressing the renderer with analytic backward passes or in
SLANG.D~\cite{SLANG.D}, to recover scene parameters from images---the inverse-rendering capability the
original proposal targeted. Second, \emph{productionising Shader Execution Reordering on Ada}:
\Cref{sec:ser} already measures a \num{1.12}--\num{1.68}$\times$ image-identical speedup from a minimal
per-bounce reorder on a rented Ada GPU; integrating it as a first-class, hardware-detected path with a
per-scene-tuned coherence hint is the natural follow-through for Ada-class deployment. Third, \emph{recompile-free capacity}: the per-ray buffers are compile-time constants whose size the
offline estimator already predicts per asset (\Cref{tab:overlap}). Removing the fixed buffer
\emph{entirely} was tried and abandoned---the cap-free streaming autopsy (\Cref{sec:autopsies}) was
bit-exact but slower. An untried middle path keeps a \emph{small, fixed} buffer but, when it fills,
re-issues the traversal from the last processed distance to gather the next chunk and repeats: this pays
an extra re-traversal only on the rare overflowing ray, keeps the single-trace fast path for the common
case, and removes the per-asset recompile without re-introducing the cap-free design's whole-render
cost. Path guiding, deferred here for want of a cheap
predictive test (\Cref{sec:autopsies}), would also merit a revisit with a full offline oracle.

\section{Closing}
\label{sec:closing}

The renderer demonstrates that ray-traced Gaussian kernel-mixture volumes can be rendered both
\emph{correctly} and \emph{efficiently} by leaning on the one property the representation affords---an
analytic, invertible per-primitive optical depth---to build a sort-free, march-free, root-find-free path
tracer. Just as usefully, the negative-results ledger of \Cref{ch:optimization} maps where further
effort on this class of renderer does and does not pay off: it is megakernel-shaped, latency-bound, and
resistant to the restructurings that help more conventional path tracers. That map is offered as a
contribution in its own right.
